% GENERATED FILE. Edit canonical lessons, then run npm run generate:books.
% canonical-source-hash: fnv1a64:c8de1c0b71dfcd0f
% canonical-lessons: RU-C134-to-tell, RU-C134-to-go-for-a-walk, RU-C134-to-tidy-up, RU-C134-to-find, RU-C134-to-look-for

\chapter{To tell, To go for a walk, To tidy up, To find, To look for}
\label{ch:ru-to-tell-to-go-for-a-walk-to-tidy-up-to-find-to-look-for}
\hlchaptermodality{134}

\begin{chapteropening}
\textbf{By the end of this chapter:} \emph{I can say to tell, to go for a walk, to tidy up, to find and to look for.}

Complete the last lesson of chapter 134: I can say to tell, to go for a walk, to tidy up, to find and to look for.
\end{chapteropening}

\section[\texorpdfstring{\ru{рассказывать} (rasskázyvat)}{рассказывать (rasskázyvat)}]{\ru{рассказывать} (rasskázyvat) — to tell (a story)}

\label{lesson:RU-C134-to-tell}



\begin{quote}
Before the new one: say the Russian for to rent; to take off, then the Russian for to repair.
\end{quote}

\subsection*{You'll want to know: \ru{рассказывать}}
\textbf{\ru{рассказывать}} (\emph{rasskázyvat}) — \textquotedblleft{}to tell (a story)\textquotedblright{}.

\begin{cousinweb}
\textbf{\ru{рассказывать}} is \emph{\ru{рас}-}, \textquotedblleft{}out, through\textquotedblright{}, with \emph{-\ru{сказывать}}, from \emph{\ru{сказать}}, \textquotedblleft{}to say\textquotedblright{}, which Vasmer derives from \emph{\ru{казать}}, \textquotedblleft{}to show\textquotedblright{}: telling is showing in words.
\end{cousinweb}

\begin{grammarlens}[title={What you've built}]
One more word for this chapter.
\end{grammarlens}

\subsection*{Your turn}
\begin{itemize}
\raggedright
  \item \emph{Say it:} \emph{rasskázyvat}
  \item \emph{Say it:} \emph{rasskázyvat}, once more
  \item \emph{Say it:} say \emph{rasskázyvat}
  \item \emph{Recall:} say the Russian for to rent; to take off, then the Russian for to repair, then say \emph{rasskázyvat} again
\end{itemize}

\begin{tcolorbox}[breakable,colback=teal!4,colframe=teal!35!black,title={Before you move on}]
What does \ru{рассказывать} mean? (\textquotedblleft{}To tell (a story)\textquotedblright{}.) Say it once more.
\end{tcolorbox}
\section[\texorpdfstring{\ru{гулять} (gulyát)}{гулять (gulyát)}]{\ru{гулять} (gulyát) — to go for a walk}

\label{lesson:RU-C134-to-go-for-a-walk}



\begin{quote}
Before the new one: say the Russian for to repair, then the Russian for to tell.
\end{quote}

\subsection*{You'll want to know: \ru{гулять}}
\textbf{\ru{гулять}} (\emph{gulyát}) — \textquotedblleft{}to go for a walk\textquotedblright{}.

\begin{cousinweb}
\textbf{\ru{гулять}} — \textquotedblleft{}a difficult word\textquotedblright{}, says Vasmer. Berneker and Brückner tied it to \emph{\ru{гул}}, \textquotedblleft{}a hum\textquotedblright{}, Potebnya and Sobolevsky to a first sense \textquotedblleft{}to drink\textquotedblright{}; each is disputed, and Trubachev's note traces the sense from \textquotedblleft{}to play ball\textquotedblright{}.
\end{cousinweb}

\begin{grammarlens}[title={What you've built}]
One more word for this chapter.
\end{grammarlens}

\subsection*{Your turn}
\begin{itemize}
\raggedright
  \item \emph{Say it:} \emph{gulyát}
  \item \emph{Say it:} \emph{gulyát}, once more
  \item \emph{Say it:} say \emph{gulyát}
  \item \emph{Recall:} say the Russian for to repair, then the Russian for to tell, then say \emph{gulyát} again
\end{itemize}

\begin{tcolorbox}[breakable,colback=teal!4,colframe=teal!35!black,title={Before you move on}]
What does \ru{гулять} mean? (\textquotedblleft{}To go for a walk\textquotedblright{}.) Say it once more.
\end{tcolorbox}
\section[\texorpdfstring{\ru{убирать} (ubirát)}{убирать (ubirát)}]{\ru{убирать} (ubirát) — to tidy up}

\label{lesson:RU-C134-to-tidy-up}



\begin{quote}
Before the new one: say the Russian for to tell, then the Russian for to go for a walk.
\end{quote}

\subsection*{You'll want to know: \ru{убирать}}
\textbf{\ru{убирать}} (\emph{ubirát}) — \textquotedblleft{}to tidy up\textquotedblright{}.

\begin{cousinweb}
\textbf{\ru{убирать}} is \emph{\ru{у}-}, \textquotedblleft{}away\textquotedblright{}, with \emph{-\ru{бирать}}, from \emph{\ru{брать}}, \textquotedblleft{}to take\textquotedblright{}: to take things away, to clear up.
\end{cousinweb}

\begin{grammarlens}[title={What you've built}]
One more word for this chapter.
\end{grammarlens}

\subsection*{Your turn}
\begin{itemize}
\raggedright
  \item \emph{Say it:} \emph{ubirát}
  \item \emph{Say it:} \emph{ubirát}, once more
  \item \emph{Say it:} say \emph{ubirát}
  \item \emph{Recall:} say the Russian for to tell, then the Russian for to go for a walk, then say \emph{ubirát} again
\end{itemize}

\begin{tcolorbox}[breakable,colback=teal!4,colframe=teal!35!black,title={Before you move on}]
What does \ru{убирать} mean? (\textquotedblleft{}To tidy up\textquotedblright{}.) Say it once more.
\end{tcolorbox}
\section[\texorpdfstring{\ru{находить} (nakhodít)}{находить (nakhodít)}]{\ru{находить} (nakhodít) — to find}

\label{lesson:RU-C134-to-find}



\begin{quote}
Before the new one: say the Russian for to go for a walk, then the Russian for to tidy up.
\end{quote}

\subsection*{You'll want to know: \ru{находить}}
\textbf{\ru{находить}} (\emph{nakhodít}) — \textquotedblleft{}to find\textquotedblright{}.

\begin{cousinweb}
\textbf{\ru{находить}} is \emph{\ru{на}-}, \textquotedblleft{}onto\textquotedblright{}, with \emph{\ru{ходить}}, \textquotedblleft{}to walk\textquotedblright{}: to find is to come upon a thing. Vasmer ties \emph{\ru{ход}} to the Indo-European root \emph{*sed-}, \textquotedblleft{}to sit\textquotedblright{}, its \emph{\ru{х}} grown from an \emph{s} after prefixes.
\end{cousinweb}

\begin{grammarlens}[title={What you've built}]
One more word for this chapter.
\end{grammarlens}

\subsection*{Your turn}
\begin{itemize}
\raggedright
  \item \emph{Say it:} \emph{nakhodít}
  \item \emph{Say it:} \emph{nakhodít}, once more
  \item \emph{Say it:} say \emph{nakhodít}
  \item \emph{Recall:} say the Russian for to go for a walk, then the Russian for to tidy up, then say \emph{nakhodít} again
\end{itemize}

\begin{tcolorbox}[breakable,colback=teal!4,colframe=teal!35!black,title={Before you move on}]
What does \ru{находить} mean? (\textquotedblleft{}To find\textquotedblright{}.) Say it once more.
\end{tcolorbox}
\section[\texorpdfstring{\ru{искать} (iskát)}{искать (iskát)}]{\ru{искать} (iskát) — to look for}

\label{lesson:RU-C134-to-look-for}



\begin{quote}
Before the new one: say the Russian for to tidy up, then the Russian for to find.
\end{quote}

\subsection*{You'll want to know: \ru{искать}}
\textbf{\ru{искать}} (\emph{iskát}) — \textquotedblleft{}to look for\textquotedblright{}.

\begin{cousinweb}
\textbf{\ru{искать}} $\leftarrow$ Proto-Slavic \emph{*iskati} — inherited. Vasmer relates it to Lithuanian \emph{ieškóti}, \textquotedblleft{}to seek\textquotedblright{}, Sanskrit \emph{iccháti}, \textquotedblleft{}seeks, wishes\textquotedblright{}, and Old English \emph{āscian}, so English \emph{ask} is a cousin.
\end{cousinweb}

\begin{grammarlens}[title={What you've built}]
One more word for this chapter.
\end{grammarlens}

\subsection*{Your turn}
\begin{itemize}
\raggedright
  \item \emph{Say it:} \emph{iskát}
  \item \emph{Say it:} \emph{iskát}, once more
  \item \emph{Say it:} say \emph{iskát}
  \item \emph{Recall:} say the Russian for to tidy up, then the Russian for to find, then say \emph{iskát} again
\end{itemize}

\begin{tcolorbox}[breakable,colback=teal!4,colframe=teal!35!black,title={Before you move on}]
What does \ru{искать} mean? (\textquotedblleft{}To look for\textquotedblright{}.) Say it once more.
\end{tcolorbox}
